% III-B Reasoning and Planning --- paste-ready subsection.
% Citations use only keys from the current thebibliography.
% \ref{sec:detect}, \ref{sec:commit}, \ref{fig:e2e} must exist in the manuscript.

\subsection{Reasoning and Planning Layer (Reason)}
\label{sec:reason}

Section~\ref{sec:detect} closes detection-without-centralization and forwards a structured record, not a narrative: the predicted class $\hat{y}_i$, the detector confidence $\mathrm{conf}_i$, the dominant evidence group $p_i^{\star}$ obtained by summing KernelSHAP over each $64$-dimensional fusion block~\cite{wu2023falcon}, the three domain shares, and the top local features by $\lvert\phi\rvert$. That record is still evidence. Knowing the event is a PortScan whose Perimeter block is largest does not name an authorized control, does not say whether Access or Endpoint should move second, and does not bind those choices to a catalog an execution gate can check. An unconstrained language model can invent a fluent mitigation that no playbook allows, or can tell a host-only story for a perimeter-dominated scan~\cite{xi2023,kshetri2025agentic,lazer2026}. We therefore treat Reasoning as a separate layer whose job is to convert the Detection product into one machine-checkable Mitigation Plans object. The layer does not authorize execution and does not write to the ledger; those failures belong to Commit and Apply.

The conversion is a fixed pipeline. After decode, nothing is free-form: the retrieval query is a filled template, ranking is a seven-step recall-then-precision cascade, the prompt is a slot-filled string, and the only accepted response is one JSON object whose action names must already appear in the per-attack catalog $W[\hat{y}_i]$ that Commit will later enforce on-chain~\cite{li2025,jan2025}.
\begin{equation}
\underbrace{(\hat{y}_i,\, p_i^{\star},\, \mathrm{conf}_i,\, \mathrm{SHAP}_i)}_{\text{III-A output}}
\;\xrightarrow{\mathrm{query}}\;
q_i
\;\xrightarrow{\mathrm{retrieve+rank}}\;
\mathcal{R}_i
\;\xrightarrow{\mathrm{prompt}}\;
\pi_i
\;\xrightarrow{\mathrm{JSON}}\;
\underbrace{\mathrm{plan}_i}_{\text{III-C payload } P}.
\label{eq:reason-pipeline}
\end{equation}

\paragraph{From Detection into a Query.}
Domain names in every later string are forced to exactly Access~/~ISP, Perimeter~/~IDS, or Endpoint~/~EDR---the same three evidence groups used to split the 88 CICFlowMeter columns~\cite{sharafaldin2018} and to aggregate SHAP (Fig.~\ref{fig:e2e}). $\hat{y}_i$ selects the allowed-action list $W[\hat{y}_i]$ and the preferred-response domains for that class. $\mathrm{conf}_i$ is written into both the query and the task block so intensity (\texttt{Immediate} $\ldots$ \texttt{Low}) tracks the detector rather than the model's tone. $p_i^{\star}$ and the three block shares become the sentence ``Dominant network domain: $\ldots$ (contribution: $x\%$)''; top-3 features per domain become retrieval keywords, and only the top-5 features by $\lvert\phi\rvert$ are admitted into the evidence JSON the model may cite. Retrieval does not ask the language model to invent a search string. The operational query is a single template filled from that record---class, confidence, dominant domain, and the three keyword lists---asking for SOC / NIST / CIS / ATT\&CK controls appropriate to that confidence. Optional rephrases exist; the evaluated path uses only the template, so two identical Detection records produce the same $q_i$. That determinism is required by the next layer: Commit hashes a plan that must be reproducible from the same evidence~\cite{zhang2024exploit}.

\paragraph{Policy Index.}
The corpus is closed---NIST SP~800-53~\cite{nist80053}, NIST SP~800-207~\cite{nist800207}, CIS Controls~\cite{cisv8}, and MITRE ATT\&CK~\cite{mitre_attack}---not an open web~\cite{shajarian2026,fu2026}. Documents are stored parent/child. A parent is a semantic section (the text the planner will read). A child is a short passage used only for nearest-neighbor search (size $384$, overlap $96$, embedding \texttt{all-MiniLM-L6-v2}), tagged with $(\mathrm{parent\_id},\,\mathrm{child\_index},\,\mathrm{source})$. Ranking operates on children; the prompt receives expanded parents. A child is precise enough to retrieve; a parent is long enough to be a control rather than a fragment that lost its if/then scope.

\paragraph{Retrieve-then-Rank Cascade.}
Dense search alone is the wrong ranking for this job. A bi-encoder returns topical neighbors and, left unchecked, fills the prompt with near-duplicates from one source while dropping a short CIS or ATT\&CK passage that is more actionable. Each step exists to close a failure of the previous one.
\begin{enumerate}
\item \emph{Dense retrieve (recall).} FAISS returns an oversampled neighbor list ($\min(300,\,20\times 6)$), keeps $20$ children, converts distance $d$ to similarity $1/(1+d)$, and balances hits across source files so one document cannot occupy the entire top-$20$. The bi-encoder is cheap and high-recall; it is not a relevance judge.
\item \emph{Merge and dedupe.} Multi-query children are unioned by $(\mathrm{parent\_id},\,\mathrm{child\_index})$ (else source, title, and a short prefix). The higher similarity wins. Without this step, overlap would inject the same passage twice and starve the later budget.
\item \emph{MMR (diversity).} Maximum Marginal Relevance with $\lambda=0.5$ keeps up to $60$ children,
\begin{equation}
\mathrm{MMR}(c)=\lambda\cdot\mathrm{sim}(q,c)-(1-\lambda)\cdot\max_{c'\in S}\mathrm{sim}(c,c').
\label{eq:mmr}
\end{equation}
After merge the pool is still redundant; without MMR the cross-encoder would re-score the same control and NIST, CIS, and ATT\&CK would not coexist in the prompt.
\item \emph{Cross-encoder (precision).} Each surviving $(q,\mathrm{passage})$ pair is scored jointly by \texttt{cross-encoder/ms-marco-MiniLM-L-6-v2} (passage truncated at $4000$ characters) and sorted. Vector similarity is a topical neighborhood, not ``this passage is a mitigation for \emph{this} PortScan at the perimeter.'' The cross-encoder is the relevance model; it is too expensive to run on the whole index, which is why steps 1--3 exist.
\item \emph{Rerank cut.} The top $20$ children by cross-encoder score are kept.
\item \emph{Parent expansion.} Each kept child is replaced by its parent section (capped at $12\,000$ characters); duplicate parents collapse to the higher child score. The planner must quote a control, not a $384$-character shard.
\item \emph{Prompt budget.} The top five sections (up to ten) are emitted as numbered $[k]$ title + body + source. If the list is empty, the slot is the explicit sentence that no relevant documents were found, and the model is forbidden to invent policy.
\end{enumerate}
Retrieval returns guidance and allowable primitives. It does not invent action names. The closed list $W[\hat{y}_i]$ is supplied in a different prompt slot and is the only legal vocabulary for \texttt{action}---the same catalog the Blockchain Governance Layer will store as an unpromptable whitelist~\cite{jan2025,karim2025}.

\paragraph{Prompt Construction.}
The user message is one template whose slots are filled in software before any decode~\cite{xi2023,li2025}. The model never sees raw telemetry, never sees the $192$-dimensional fusion vector, and never sees action names outside $W[\hat{y}_i]$. The slots are: (i)~role and the three exact domain strings; (ii)~prediction summary ($\hat{y}_i$, $\mathrm{conf}_i$); (iii)~the dominance sentence produced from $p_i^{\star}$; (iv)~the Detection evidence JSON, top-$5$ features only; (v)~$W[\hat{y}_i]$ plus preferred domains and typical evidence cues for that class; (vi)~a per-domain card (role, top-$3$ features, actions that domain is allowed to execute---the intersection of class whitelist and domain capability); (vii)~the numbered RAG sections, or the empty-KB sentence. The task block, also template text, requires the model to interpret $\hat{y}_i$ on all three domains, name which domain must act first, select actions only from the list, assign each a network tier, cite evidence in per-action reasoning, scale execution priority with $\mathrm{conf}_i$, and return empty lists rather than a made-up control if nothing fits. Decoding uses GPT-4o-mini at temperature $T=0.3$ with the instruction to return only valid JSON; the first ``\{'' through the last ``\}'' is parsed and surrounding prose is discarded.

\paragraph{Response Schema.}
The only accepted output is one object $\mathrm{plan}_i$ (Table~\ref{tab:plan-schema}). The field \texttt{action} is copied from $W[\hat{y}_i]$, not paraphrased (\texttt{limit rate} remains \texttt{limit rate}). \texttt{network\_tier} is exactly one of the three domain strings. Primary actions prefer $p_i^{\star}$ unless the evidence contradicts it. Empty RAG forbids invented policy citations. A PortScan whose Perimeter share dominates is therefore planned as, for example, \texttt{tarpit scan} and \texttt{harden ports} at Perimeter~/~IDS with a supporting \texttt{update ACL} at Access~/~ISP---names that are legal for PORTSCAN and illegal for BENIGN~\cite{sharafaldin2018}. That legality is not left to the model: Apply will reject any name that is absent from $W[\hat{y}_i]$ or absent from this object.

\begin{table}[t]
\caption{Mitigation Plans object $\mathrm{plan}_i$ accepted from the Reasoning layer. Action names must occur in $W[\hat{y}_i]$; $\mathtt{network\_tier}$ must be one of the three enterprise domains.}
\label{tab:plan-schema}
\centering
\begin{tabular}{p{0.30\columnwidth}p{0.62\columnwidth}}
\hline
Field & Constraint \\
\hline
\texttt{threat\_level} & $\{\mathrm{Critical},\mathrm{High},\mathrm{Medium},\mathrm{Low}\}$ \\
\texttt{all\_actions} & Union of primary and supporting names \\
\texttt{primary\_actions}[] & $\{\mathtt{action},\,\mathtt{network\_tier},\,\mathtt{party\_evidence\_type},\,\mathtt{reasoning}\}$ \\
\texttt{supporting\_actions}[] & Same shape as primary units \\
\texttt{overall\_reasoning} & Why this Access~/~Perimeter~/~Endpoint order \\
\texttt{execution\_priority} & $\{\mathrm{Immediate},\mathrm{High},\mathrm{Standard},\mathrm{Low}\}$ \\
\texttt{knowledge\_sources\_used} & Catalog and/or retrieved-section tags \\
\hline
\end{tabular}
\end{table}

\paragraph{Handoff to Blockchain Governance.}
$\mathrm{plan}_i$ is this layer's only output. It is the off-chain payload $P$ that Section~\ref{sec:commit} will canonicalize and commit as $c=\mathrm{SHA256}(P)$~\cite{zhang2024exploit,jan2025}. Reason never calls the contract and never opens a device API. This closes the evidence-to-plan failure that Section~\ref{sec:detect} left open: the planner cannot speak an action that is not in the catalog, cannot omit a domain assignment, and cannot hide the Detection fields it used. It does not close the trust failure that follows generation~\cite{chhabra2026,pendli2025}. A fluent, catalog-legal action copied from another cycle, a unit whose network tier is rewritten after decode, or a payload whose bytes change before dispatch will still look like a valid plan if the next layer only stores text. Binding $W[\hat{y}_i]$ on-chain, writing $c$ once, and refusing any apply that does not re-hash to $c$ is the failure the Blockchain Governance Layer exists to close.
